% LPAR-26 short presentation paper (Kalpa Publications).
% NOTE: copy `easychair.cls` and `lstlean.tex` from the main paper's directory next to this file.
\documentclass{easychair}

\usepackage{amsmath}
\usepackage{amssymb}
\usepackage{booktabs}
\usepackage{listings}
\usepackage{xcolor}
\usepackage{csquotes}
\usepackage{tikz}
\usepackage{cleveref}
\usepackage{fontawesome}

\newtheorem{thm}{Theorem}[section]
\newtheorem{theorem}[thm]{Theorem}
\newtheorem{definition}{Definition}[section]
\newtheorem{lemma}[thm]{Lemma}

\definecolor{academicpurple}{rgb}{0.4,0.0,0.35}
\newcommand{\secref}[1]{\hyperref[#1]{\textcolor{academicpurple}{\S\ref*{#1}}}}
\newcommand{\citeref}[1]{\textcolor{academicpurple}{\cite{#1}}}
\newcommand{\thmref}[1]{\textcolor{academicpurple}{\hyperref[#1]{Theorem~\ref*{#1}}}}
\newcommand{\extlink}{~\ensuremath{{}^\text{\faExternalLink}}}

\definecolor{keywordcolor}{rgb}{0.7, 0.1, 0.1}
\definecolor{tacticcolor}{rgb}{0.0, 0.1, 0.6}
\definecolor{commentcolor}{rgb}{0.4, 0.4, 0.4}
\definecolor{symbolcolor}{rgb}{0.0, 0.1, 0.6}
\definecolor{sortcolor}{rgb}{0.1, 0.5, 0.1}
\definecolor{attributecolor}{rgb}{0.7, 0.1, 0.1}
\definecolor{codebg}{rgb}{0.96,0.96,0.99}
\definecolor{codeframe}{rgb}{0.75,0.70,0.80}
\def\lstlanguagefiles{lstlean.tex}
\lstset{language=lean,
  basicstyle=\fontsize{8.4}{9.8}\selectfont\ttfamily,
  backgroundcolor=\color{codebg},
  frame=single, framerule=0.5pt, rulecolor=\color{codeframe},
  framesep=4pt, xleftmargin=6pt, xrightmargin=6pt,
  aboveskip=4pt, belowskip=4pt}

\title{A Certified Encoding of $n$-Queens Completion\\ as a Hopfield/Boltzmann Network}

\titlerunning{A Certified Encoding of $n$-Queens Completion}

\author{
Matteo Cipollina\inst{1}
\and
Michail Karatarakis\inst{2}
\and
Freek Wiedijk\inst{2}
}

\institute{
  Axiomatic AI, Barcelona, Spain\\
  \email{matteo.cipollina@yahoo.it}
\and
  Radboud University, Nijmegen, The Netherlands\\
  \email{michail.karatarakis@ru.nl, freek@cs.ru.nl}\\
}

\authorrunning{Cipollina, Karatarakis and Wiedijk}

\begin{document}
\maketitle

\begin{abstract}
In our LPAR-26 paper we formalized Hopfield networks and Boltzmann machines in Lean~4, and remarked
in passing that minimizing a Hopfield energy is the same thing as quadratic unconstrained binary
optimization. This paper works out that remark. We formalize a class of optimization problems in the
form $\lVert A x - b\rVert^2$, prove that each one is a Boltzmann network whose lowest-energy states
are exactly its solutions, and use this to encode $n$-Queens Completion: given a chessboard with
some queens already placed, can the placement be extended so that no two queens attack each other?
We prove the encoding correct in both directions. The second direction is the useful one: it means
that when the solver reports that no solution exists, that report is a theorem about the chessboard
and not a failure of the search. We give a $7 \times 7$ board with two given queens that do not
attack each other and that admits no completion, and we prove that it admits none.
\end{abstract}

\section{Introduction}\label{sec:intro}

Our LPAR-26 paper~\citeref{hnbm} formalizes Hopfield networks and Boltzmann machines in Lean~4. In
its introduction we note that minimizing the energy of a Hopfield network is, after an affine change
of variables, the problem known as quadratic unconstrained binary optimization
(QUBO)~\citeref{lucas2014ising}, and that a machine-checked account of the dynamics would give a
specification against which algorithms running on annealing hardware could be verified. We did not
carry that out. This paper carries it out for one problem.

The plan is as follows. We fix a class of constraint problems and give them a numerical objective
function which is zero exactly on the solutions (\secref{sec:qubo}). We show that the objective is
the energy of a Boltzmann network of~\citeref{hnbm}, so that solving the constraint problem and
minimizing the energy of the network are the same task (\secref{sec:net}). We then encode a
particular \textsf{NP}-complete problem in this form and prove that the encoding is faithful
(\secref{sec:queens}--\secref{sec:correct}). Finally we run the neurodynamic solver of Li and
Wang~\citeref{liwang2022} on it and report what happens (\secref{sec:solver}--\secref{sec:meas}).

We should say at the outset what is \emph{not} being claimed. We do not claim that the resulting
solver is fast; a short backtracking program beats it on every instance we tried, and we give the
numbers. We do not claim any new algorithm; the solver is Li and Wang's, used unchanged. What is new
is that the translation from the chessboard to the network is proved correct, in both directions,
and that the solver's negative answers therefore mean something.

\section{Constraint problems as sums of squares}\label{sec:qubo}

We begin with the class of problems we can handle, and we begin concretely.

Suppose we have $n$ Boolean unknowns $x_1,\dots,x_n$, each either $0$ or $1$, and $m$ constraints,
each of which says that a certain subset of the unknowns must contain exactly one $1$. Write $A$ for
the $m \times n$ matrix whose entry $A_{r,u}$ is $1$ if unknown $u$ occurs in constraint $r$ and $0$
otherwise, and write $b$ for the column vector of $m$ ones. Then the constraint \enquote{row $r$
contains exactly one $1$} says precisely that the $r$-th entry of $Ax$ equals the $r$-th entry
of~$b$. Solving all the constraints at once is solving the linear system $Ax = b$ over $\{0,1\}$.

Rather than solve the system, we measure how badly a given $x$ fails to solve it. Define
\begin{equation}\label{eq:p}
  p(x) \;=\; \lVert Ax - b\rVert^2 \;=\; \sum_{r=1}^{m}\Bigl(\ \sum_{u \in r} x_u \;-\; b_r\ \Bigr)^{\!2}.
\end{equation}
Each summand is the square of an integer, so $p(x) \ge 0$ always, and $p(x) = 0$ if and only if
every constraint is met exactly. We call $x$ \emph{feasible} when $p(x) = 0$. Minimizing $p$ is
therefore the same as solving the constraints, and $p$ is a quadratic polynomial in Boolean
variables, which is what makes the connection to Hopfield energies in \secref{sec:net} possible.

Two observations about \eqref{eq:p} will be used repeatedly, and both are easy.

\begin{lemma}[integrality]\label{lem:int}
$p(x)$ is a non-negative integer for every $x$, so if $p(x) \neq 0$ then $p(x) \ge 1$.
\end{lemma}

There is no state whose objective value is a small positive fraction. A search either lands exactly
on a solution or misses by at least a whole unit. We will see this in the measurements of
\secref{sec:meas}, where every failed run on an unsolvable instance stops at exactly $1$.

\begin{lemma}[at-most-one by a slack variable]\label{lem:slack}
Let $S$ be a set of unknowns and let $s$ be a fresh Boolean unknown. Then the constraint
$\sum_{u \in S} x_u + s = 1$ is satisfiable if and only if $\sum_{u \in S} x_u \le 1$, and in that
case $s$ is determined: $s = 1 - \sum_{u\in S} x_u$.
\end{lemma}

This is the only trick we need in order to express \enquote{at most one} in a language that only has
\enquote{exactly one}. If nothing in $S$ is set, $s$ takes up the slack and is set instead. If
exactly one thing in $S$ is set, $s$ is $0$. If two or more are set, no value of $s$ works, the row
contributes at least $1$ to \eqref{eq:p}, and the assignment is penalized. Note that a single slack
suffices however large $S$ is; we do not need one slack per pair.

\medskip\noindent\textbf{Which problems fit.} The form \eqref{eq:p} accepts exactly the
\emph{feasibility} problems whose constraints are equalities, possibly after adding slack variables
as in \cref{lem:slack}. Consulting the standard catalogue~\citeref{lucas2014ising}: exact cover
(\S4.1 there) fits with no slack variables at all; graph colouring (\S6.1) fits after one slack per
edge and colour, which gives a problem equivalent to, though not literally identical with, the
Hamiltonian given there, whose conflict term is a bare product $x_u x_v$ rather than a squared
residual. Optimization problems such as vertex cover (\S4.3) and the travelling salesman (\S7.2) do
not fit, because they ask for a \emph{best} solution rather than for \emph{any} solution, and
\eqref{eq:p} has no way to express a preference among feasible points. We say this because it
delimits the method: what follows applies to a restricted class, and it is better to be explicit
about which.

The catalogue does not include $n$-queens, in any of its sections. The encoding of
\secref{sec:enc} is therefore not a transcription of a published Hamiltonian but one we
constructed, which is a further reason to prove it correct rather than to assume it.

\medskip\noindent\textbf{In Lean.} The data above is a structure we call \lstinline{Problem}. It
stores the number of unknowns and rows, the matrix $A$ column-wise (for each unknown, the list of
rows it occurs in), the vector $b$, and two derived quantities that \secref{sec:net} needs. Nothing
in it mentions any particular application; graph colouring, exact cover, Sudoku and the queens
problem of this paper are all values of this one structure.

\section{The problem: $n$-Queens Completion}\label{sec:queens}

A queen on a chessboard attacks along its row, its column, and both diagonals. The classical
$n$-queens problem asks for a placement of $N$ queens on an $N \times N$ board, no two attacking. It
is a familiar puzzle, and for our purposes it has a fatal defect: it is solvable for \emph{every}
$N \ge 4$, by a closed-form construction. A theorem of the shape \enquote{the encoding has a
solution if and only if the board does} would therefore be true but empty, since both sides are
always true.

We use instead the completion variant.

\begin{definition}[$n$-Queens Completion]\label{def:completion}
An instance is a board size $N$ together with a finite list of \emph{givens}, that is, squares
$(i,j)$ on which a queen has already been placed. The instance is a \textsf{yes} instance when there
is a placement of $N$ mutually non-attacking queens, one in each row, that includes every given.
\end{definition}

This variant is \textsf{NP}-complete, and counting its solutions is $\#\textsf{P}$-complete
\citeref{gent2017}. It is also much better behaved as a test of an encoding, because whether an
instance is solvable genuinely depends on the givens, and one cannot tell by inspection. Our corpus
contains a $7 \times 7$ board whose two given queens are on different rows, different columns and
different diagonals --- they do not attack each other, so nothing is locally wrong --- and which
nevertheless admits no completion at all.

It is worth seeing how quickly this becomes non-obvious. A queen in the corner square $(0,0)$ is
fatal on a $4 \times 4$ board and on a $6 \times 6$ board, and harmless on every other size we
checked: counting the completions of $\langle n, \{(0,0)\}\rangle$ for $n = 4,\dots,10$ gives
$0,\,2,\,0,\,4,\,4,\,28,\,64$.

We represent a board as an array $q$ of $N$ natural numbers, where $q_i$ is the column of the queen
in row $i$. Placing one queen per row is then built into the representation, and the remaining
conditions are checked by a Boolean function \lstinline{isQueens}: that $q$ has length $N$, that
every entry is less than $N$, that no two rows share a column or a diagonal, and that every given is
present.

\medskip\noindent\textbf{Diagonals without subtraction.} Two queens at $(i,q_i)$ and $(k,q_k)$ share
a diagonal when $i - q_i = k - q_k$, and an anti-diagonal when $i + q_i = k + q_k$. The first of
these involves a subtraction, and in Lean subtraction on $\mathbb{N}$ is truncated: $2 - 5$ is $0$,
not $-3$. Writing the condition that way would force a case distinction, or a move to $\mathbb{Z}$,
at every use. We therefore rearrange it to $i + q_k = k + q_i$, which is equivalent over the
integers and involves only addition. The same rearrangement is used throughout the encoding below.
It is a small point, but it removes a great deal of case analysis, and it is the kind of thing that
is invisible on paper and unavoidable in a formalization.

\section{The encoding}\label{sec:enc}

We now write $n$-Queens Completion in the form of \secref{sec:qubo}. Take one Boolean unknown
$x_{i,j}$ for each square, to be read as \enquote{there is a queen on square $(i,j)$}. The
conditions of \cref{def:completion} become the following rows.

\begin{center}\small
\begin{tabular}{@{}llc@{}}
\toprule
condition & row & how many\\
\midrule
row $i$ holds exactly one queen & $\sum_j x_{i,j} = 1$ & $N$\\
column $j$ holds exactly one queen & $\sum_i x_{i,j} = 1$ & $N$\\
diagonal $d$ holds at most one & $\sum_{(i,j) \in d} x_{i,j} + s_d = 1$ & $2N-3$\\
anti-diagonal $a$ holds at most one & $\sum_{(i,j) \in a} x_{i,j} + s_a = 1$ & $2N-3$\\
the given at $(i,j)$ is present & $x_{i,j} = 1$ & one per given\\
\bottomrule
\end{tabular}
\end{center}

The first two families are equalities already. The diagonal families are at-most-one conditions, and
we express them using \cref{lem:slack}, with one fresh slack variable per diagonal. The givens are
rows containing a single unknown, whose only solution is to set it; this is why we need no separate
mechanism for fixing variables in advance.

Only diagonals of length at least two are given rows. A diagonal consisting of a single square
cannot contain two queens, so the condition on it is vacuous, and a row for it would contain one
unknown and one slack and would say nothing. On an $N \times N$ board there are $2N-1$ diagonals in
each direction, of which the two extreme ones are single corner squares; hence $2N-3$ rows in each
direction. Counting up,
\[
  \#\text{unknowns} = N^2 + 2(2N-3) = N^2 + 4N - 6,
  \qquad
  \#\text{rows} = 2N + 2(2N-3) + \#\text{givens} = 6N - 6 + \#\text{givens}.
\]
For the ordinary $8 \times 8$ board this is $90$ unknowns and $42$ rows.

\medskip\noindent\textbf{A worked instance.} Take $N = 4$ with one given queen at the corner
$(0,0)$. There are $16$ square unknowns and $2 \times 5 = 10$ slacks, so $26$ in all; and
$4 + 4 + 5 + 5 = 18$ rows, plus one more for the given, so $19$. This instance has no completion:
the $4 \times 4$ board has exactly two solutions, $(1,3,0,2)$ and $(2,0,3,1)$, and neither places a
queen in column $0$ of row $0$. We return to it in \secref{sec:correct}.

\medskip\noindent\textbf{The degrees are uneven, and this matters.} In \secref{sec:net} the network
attached to a \lstinline{Problem} has a threshold $\theta_u$ at each unknown $u$, given by
$\theta_u = \tfrac{1}{2}\deg(u) - \sum_{r \ni u} b_r$, where $\deg(u)$ is the number of rows
containing $u$. In our encoding every $b_r$ is $1$, so the sum is just $\deg(u)$ and
$\theta_u = -\tfrac{1}{2}\deg(u)$.

Counting degrees: a square in the middle of the board occurs in four rows --- its board row, its
column, its diagonal and its anti-diagonal. Each of the four corner squares lies on a diagonal of
length one, which carries no row, so it occurs in only three. A slack occurs in exactly one. Hence
three distinct degrees, $4$, $3$ and $1$, two of them odd, and correspondingly
$\theta_u \in \{-2,\,-1.5,\,-0.5\}$.

Half-integer thresholds are unproblematic for the network of~\citeref{hnbm}, whose $\theta$ is
real-valued. They are a problem for us only because we want the solver's inner loop to run in exact
integer arithmetic, and so store $\theta$ in an integer array. An earlier version of the library did
store $\theta$ itself, and kept every degree even by writing each constraint twice, which doubled the
number of rows, the objective value and the work per sweep. That workaround was removed by storing
$2\theta_u$, which is an integer whatever the degrees are, and halving it at the single point where
the real-valued energy is defined. Sudoku, where every unknown has degree exactly four, would never
have prompted this. We report the queens degrees because they exercise the resulting degree-free
algebra harder than any other instance we have: three different values, two of them odd, arising
from the geometry of the board rather than from a synthetic test case.

\section{The encoding is faithful}\label{sec:correct}

Two functions connect boards and assignments. Given a board $q$, the function \lstinline{encode}
produces the assignment that sets $x_{i,j}$ exactly when $q_i = j$, and sets each slack to the value
forced by \cref{lem:slack}. Given an assignment $x$, the function \lstinline{decode} reads off, for
each row $i$, the first column $j$ with $x_{i,j}$ set.

\begin{thm}\label{thm:iff}
Suppose every given queen lies on the board. Then:
\begin{enumerate}
\item[(i)] (soundness) if $p(x) = 0$ then \lstinline{isQueens (decode x)} holds;
\item[(ii)] (completeness) if \lstinline{isQueens q} holds then $p(\texttt{encode } q) = 0$;
\item[(iii)] consequently $\bigl(\exists x.\; p(x) = 0\bigr) \iff \bigl(\exists q.\;
      \texttt{isQueens } q\bigr)$.
\end{enumerate}
\end{thm}

The hypothesis is the only one. In particular there is no lower bound on $N$. We did not exclude the
degenerate boards but discharged them: when $N = 0$ there are no unknowns and no rows, so $p$ is the
empty sum and is zero, and the empty board is vacuously a completion, so both sides of (iii) hold.
The one case that must be excluded is a given lying off the board, because the row for such a given
would contain no unknown at all and could never be satisfied; the hypothesis is exactly what rules
this out. We are explicit about this because a reader is entitled to suspect that a side condition
of the form \enquote{assume the board is non-empty} is concealing something.

\medskip\noindent\textbf{Why the second direction is the interesting one.} Soundness says that when
the solver reports a zero, the board it produces is a genuine solution. That is reassuring, but one
could get the same assurance by checking the answer afterwards, which is cheap. Completeness says
something a checker cannot: that if there had been a solution, the objective would have had a zero.
Contrapositively, if the objective has no zero, there is no solution. So a run that ends without
finding anything is evidence about the chessboard rather than about the search, and statements like
the following become available.

\begin{center}\small
\begin{tabular}{@{}lllc@{}}
\toprule
instance & board and givens & why it is blocked & size of the check\\
\midrule
\texttt{attacking} & $8\times8$; $(0,0),(1,1)$ & the givens share a diagonal & none needed\\
\texttt{blocked4}  & $4\times4$; $(0,0)$ & a corner queen is fatal at $N=4$ & $4^3 = 64$\\
\texttt{blocked6}  & $6\times6$; $(0,0)$ & and at $N=6$ & $6^5 = 7776$\\
\texttt{blocked7}  & $7\times7$; $(0,0),(1,6)$ & the givens do \emph{not} attack & $7^5 = 16807$\\
\bottomrule
\end{tabular}
\end{center}

Each of these is a theorem of the form $\neg\,\exists q.\ \texttt{isQueens}\ q$, that is, a statement
about chessboards, with no mention of the encoding. They are proved by exhausting the possibilities
that the givens leave open: the checker looks only at the first $N$ entries of the array, so a
hypothetical completion is one of finitely many tuples, and the givens pin some coordinates. The
enumeration is performed by the Lean kernel (\lstinline{decide +kernel}); we do not use
\lstinline{native_decide} anywhere, so no compiler or runtime is part of the trusted base. By
\thmref{thm:iff}(iii) each then upgrades to the statement that no assignment whatsoever reaches
$p = 0$ --- a claim about $2^{71}$ assignments in the case of \texttt{blocked7}, obtained from a
check over $16807$ tuples.

\section{The objective is the energy of a Boltzmann network}\label{sec:net}

We now connect \secref{sec:qubo} to~\citeref{hnbm}. A Hopfield energy has the shape
$-\tfrac{1}{2}\sum_{u \neq v} W_{uv}\,x_u x_v + \sum_u \theta_u x_u$, with $W$ symmetric and zero on
the diagonal. Expanding \eqref{eq:p} gives $p(x) = x^{\top}\!A^{\top}\!A\,x - 2b^{\top}\!Ax +
\lVert b\rVert^2$, which is quadratic but has a diagonal: the term $(A^{\top}A)_{uu}x_u^2$. Since
$x_u$ is Boolean we may replace $x_u^2$ by $x_u$, which turns that term into a linear one and lets
it be absorbed into the threshold. What is left off the diagonal is symmetric and vanishes on it.
Setting $W = -(A^{\top}A - \operatorname{diag}(A^{\top}A))$ and folding the rest into $\theta$ gives
the following.

\begin{thm}\label{thm:net}
For any \lstinline{Problem}: the energy of the associated $\{0,1\}$ Boltzmann network
of~\citeref{hnbm} at the state corresponding to $x$ equals
$\tfrac{1}{2}\bigl(p(x) - \lVert b\rVert^2\bigr)$. Hence a state has minimum energy
$-\tfrac{1}{2}\lVert b\rVert^2$ exactly when $p(x) = 0$, and by \cref{lem:int} every state that is
not feasible has energy at least $\tfrac{1}{2}$ above the minimum.
\end{thm}

The energy gap of $\tfrac{1}{2}$ is worth isolating. Because the state space is finite and the gap
does not shrink, the Boltzmann weight $e^{-E/T}$ of an infeasible state is smaller than that of a
feasible one by a factor of at least $e^{-1/2T}$, which tends to $0$ as the temperature $T$ falls.
Combining this with the ergodicity result of~\citeref{hnbm} --- which says that the random-scan Gibbs
chain has a unique stationary distribution, namely the Boltzmann law --- gives that the chain's
equilibrium distribution concentrates on the solutions as $T \to 0^{+}$. At the queens instance we
also prove the converse pairing, that the minimum-energy states are exactly the completions, in both
directions.

\medskip\noindent\textbf{Where the certification stops.} We are careful to distinguish three things.
The \emph{objective} is the network's energy, by \thmref{thm:net}. The \emph{quantity computed in
the solver's inner loop} is the network's local field: the solver maintains its own integer routine
for efficiency, and we prove that routine equal to the field defined in~\citeref{hnbm}. The
\emph{trajectory the solver follows}, however, is not covered. Equations (3) and (6)
of~\citeref{liwang2022} accumulate their input, $u \leftarrow u + (Wx - \theta)$, and this
accumulation is not the memoryless single-site update for which the convergence and detailed-balance
theorems of~\citeref{hnbm} are stated; the two agree only on the first inner iteration, when $u$ is
still zero. So the objective, the field, and the minimizers are certified, and the path taken
through the state space is not.

\section{The solver}\label{sec:solver}

The search is Algorithm~2 of Li and Wang~\citeref{liwang2022}, used unchanged. It maintains a
population of $N$ copies of a neurodynamic model. Each copy is run until it settles; the best state
found so far, the \emph{incumbent}, is recorded; the population is then re-initialized around the
incumbent by a particle-swarm rule (eq.~7 there), and a measure of how spread out the population is
(eq.~8) triggers random bit flips (eq.~9) when the copies have collapsed onto each other. The search
stops after $M$ consecutive rounds without improving the incumbent. Each copy is either the
deterministic recurrence of eq.~(3), which we call DHNm, or its stochastic counterpart eq.~(6),
which we call BMm and which accepts a flip with a probability given by a logistic function of the
local field divided by a temperature. We implement both.

\medskip\noindent\textbf{Only this stage applies here.} Li and Wang precede the search with an
Algorithm~1, which is a deterministic simplification step: it is Sudoku constraint propagation, and
on five of their ten benchmark puzzles it finishes the job by itself, so that the search never runs.
Our queens pipeline has no such stage; the encoding is built directly from the board and nothing is
assigned in advance. An analogue exists --- if some row or column has exactly one square not attacked
by a given, a queen must go there, and one can iterate --- and would shrink the encoding on heavily
constrained instances. We have not implemented it.

\medskip\noindent\textbf{The constants the source does not give.} Reproducing~\citeref{liwang2022}
requires values for the initial temperature, the cooling rate, the number of inner iterations, the
mutation probability, the diversity threshold $\mathcal{T}$, and the three particle-swarm
coefficients $c_0,c_1,c_2$. None is stated anywhere in that paper, so all had to be recovered
experimentally. Two of them turned out to matter in a way that goes beyond tuning.

The first is $c_2$, which weights the pull of the incumbent. If $c_2$ is comparable to $c_1$, every
copy is dragged onto the incumbent; once the incumbent sits on a plateau the population collapses
and the search stops exploring. Taking $c_2$ well below $c_1$ removes this.

The second is $\mathcal{T}$. Equation~(8) measures diversity as an average of Euclidean distances
divided by the number of unknowns $n$, and this quantity is at most $1/\sqrt{n}$. A fixed threshold
therefore means something different on each instance, and on a large instance it can fire on every
round. We compare $\mathcal{T}$ against the diversity scaled by $\sqrt{n}$, which makes it a number
in $[0,1]$ with the same meaning everywhere. The value we recovered is model-specific, $0.9$ for
DHNm against $0.4$ for BMm. It was fitted on Sudoku, and it transfers: on the empty $8 \times 8$
queens board DHNm fails at BMm's $0.4$ and succeeds on $9$ of $10$ seeds at its own $0.9$. A constant
recovered on one problem and confirmed on an unrelated one is on much firmer ground than a fitted
one.

\section{What happens when we run it}\label{sec:meas}

All runs use population and stall limit $20$, at most $60$ rounds, $40$ inner iterations, initial
temperature $3.0$, cooling factor $0.9$, and a fixed base seed; the generator is a
\texttt{splitmix64} and nothing reads a clock, so the table can be reproduced exactly. The
comparison is against a backtracking completion search written separately and sharing nothing with
the encoding. That baseline is itself certified: we prove that every board it returns satisfies
\lstinline{isQueens}. We do not prove the converse, that it finds a solution whenever one exists;
it is exhaustive by construction, but that is an argument about the program rather than a theorem
about it, which is why the blocked instances of \secref{sec:correct} are refuted separately.

\begin{center}\small
\begin{tabular}{@{}lrrr rr rrr r r@{}}
\toprule
& & & & \multicolumn{2}{c}{backtracking} & \multicolumn{3}{c}{search, with BMm} & DHNm & one BMm\\
\cmidrule(lr){5-6}\cmidrule(lr){7-9}\cmidrule(lr){10-10}\cmidrule(lr){11-11}
board & $N$ & givens & unknowns & ok & solutions & found & rounds & $s/k$ & $s/k$ & $s/20$\\
\midrule
Q4     & 4  & 0 & 26  & yes & 2   & yes & 1  & 10/10 & 10/10 & 7\\
Q6     & 6  & 0 & 54  & yes & 4   & yes & 4  & 9/10  & 8/10  & 0\\
Q8     & 8  & 0 & 90  & yes & 92  & yes & 3  & 10/10 & 9/10  & 0\\
Q10    & 10 & 0 & 134 & yes & 724 & no  & 31 & 0/5   & 1/5   & 0\\
Comp6  & 6  & 1 & 54  & yes & 1   & no  & 22 & 6/10  & 7/10  & 0\\
Comp8  & 8  & 3 & 90  & yes & 1   & yes & 14 & 6/10  & 3/10  & 0\\
Comp8b & 8  & 5 & 90  & yes & 1   & yes & 4  & 10/10 & 10/10 & 0\\
\bottomrule
\end{tabular}
\end{center}

The columns are, in order: the size of the board and of its encoding; whether backtracking found a
completion and whether the checker accepts the board it returned; the number of completions; whether
the search found a zero at the base seed and in how many rounds; how many of $k$ seeds succeeded;
the same for DHNm in place of BMm; and how many of $20$ single annealed runs, with no population at
all, ended on a completion.

Three things are worth drawing out.

\emph{Backtracking wins, comfortably.} It solves every solvable instance and refutes every
unsolvable one, in times too small to be worth tabulating. We report this rather than working around
it. The contribution here is that the answers are proved correct, not that they are obtained
quickly.

\emph{A single annealed machine is much weaker than the population.} It succeeds on $7$ of $20$
seeds at $N = 4$ and $5$ of $20$ at $N = 5$, and on nothing at all from $N = 6$ upwards or on any
instance with givens, whereas the population solves ten of the eleven boards in our corpus. This is
a larger gap than we see for graph colouring, and it is the clearest evidence we have that the
collaborative part of the method is doing the work.

\emph{Which of the two models one uses matters here and not elsewhere.} Under the population search
DHNm keeps pace with BMm on the empty boards but falls behind on the instances with many givens.
On our graph-colouring corpus the two are indistinguishable: DHNm attains the chromatic number on
$10$ of $10$ seeds for fourteen of the fifteen graphs and fails on the fifteenth, exactly matching
BMm. Neither corpus on its own would have shown this.

On the unsolvable instances the search finds no zero on any seed, and the best value it reaches is
exactly $1$ on every one of them, never a fraction. This is \cref{lem:int} appearing in the output.

The count of solutions is a check on the baseline rather than a result: on empty boards it returns
$1,1,0,0,2,10,4,40,92,352,724$ for $n = 0,\dots,10$, which is the classical sequence, including the
$92$ solutions of the ordinary eight-queens problem. It also exercises the counting version of the
problem, which the decision encoding alone does not reach.

\section{Watching it run}\label{sec:widget}

The development includes widgets that draw a run in the Lean infoview as a chessboard, one frame per
round of the search, with a queen on each occupied square, given queens drawn in black and placed
queens in white, and a dashed line joining every attacking pair. The line is drawn exactly when one
of the three pairwise conditions of \lstinline{isQueens} fails, which is exactly when some column or
diagonal row of $A$ contains two set unknowns; so what one watches disappearing is the residual of
\eqref{eq:p} and not an illustration of it.

One widget is interactive rather than a replay: the reader places queens by clicking on squares and
then runs the solver on the resulting partial board. This is done with Lean's remote-procedure-call
mechanism for widgets, so the code that runs is the same code the theorems are about. It requires
the file to be open in the editor and runs in the Lean interpreter, so it is a tool for reading the
development rather than a web page.

We should record one thing about these animations, since it bears on \secref{sec:meas}. The
single-machine widgets are pinned to seeds that succeed, and they anneal five times more slowly than
the default schedule; at the default schedule no seed among the first sixty succeeds on either of the
two instances with givens. The animations are chosen to end on a solved board. The honest summary of
what one machine can do is the last column of the table above.

\section{Related work}\label{sec:related}

\noindent\textbf{The problem.} Bell and Stevens survey what is known about
$n$-queens~\citeref{bell2009}. The completion question goes back to Nauck in 1850; Gent, Jefferson
and Nightingale proved it \textsf{NP}-complete and $\#\textsf{P}$-complete~\citeref{gent2017}, and
Glock, Munhá Correia and Sudakov showed that any placement of at most $n/60$ mutually non-attacking
queens can be completed while some placements of about $n/4$ cannot~\citeref{glock2022}. Their bound
is asymptotic, so it does not speak to our $7\times7$ example of \secref{sec:correct}. Counting
solutions of the unconstrained problem is a separate line of work~\citeref{simkin2021}.

\medskip\noindent\textbf{Neural approaches to $n$-queens.} These begin with the Hopfield--Tank
programme for optimization and the criticisms of it by Wilson and Pawley. Mańdziuk applied a binary
Hopfield-type network directly to $n$-queens~\citeref{mandziuk1995}; as \secref{sec:notclaim}
records, it is a better solver than ours. Mérida, Muñoz and Benítez used $n$-queens as the benchmark
for a multivalued recurrent network~\citeref{merida2001}. The solver we use is an instance of the
collaborative neurodynamic framework of Che and Wang~\citeref{che2019}, which does come with a
global convergence claim --- but for an augmented-Lagrangian neurodynamic model, not for the
momentum recurrences of eqs.~(3) and~(6) that~\citeref{liwang2022} actually iterates, and this is
the gap \secref{sec:net} describes.

\medskip\noindent\textbf{QUBO encodings, and at-most-one in particular.} Lucas's
catalogue~\citeref{lucas2014ising} and the tutorial of Glover, Kochenberger and
Du~\citeref{glover2019} are the standard references; neither treats $n$-queens. The closest work to
our \secref{sec:enc} is Codognet~\citeref{codognet2023}, who compares several QUBO encodings of the
at-most-one constraint \emph{using $n$-queens as the example domain}, on a D-Wave annealer and on a
quantum-inspired engine. That is the same constraint on the same problem, and the difference is
exactly the one this paper is about: the encodings there are compared by measurement, and none is
proved to be faithful.

\medskip\noindent\textbf{Verified encodings.} That a solver can be trustworthy while the encoding
handed to it is not is by now a recognised gap, and it has been attacked for propositional
formalisms. Codel, Avigad and Heule verify CNF encodings in Lean, including at-most-one and
at-most-$k$~\citeref{codel2023}; Szeider's PBLean imports pseudo-Boolean certificates into Lean~4
and, in the same spirit as this paper, formalizes both the constraint translation and its
correctness proof so that theorems about the original problem follow~\citeref{szeider2026}. On the
solver side, the VeriPB line of work provides certificates for pseudo-Boolean and constraint
solving~\citeref{bogaerts2022}, and Subercaseaux and Heule adopted a verified approach precisely
because their encoding had become too intricate to trust~\citeref{subercaseaux2023}. Our
contribution is the analogous step for the canonical form $\lVert Ax-b\rVert^2$ rather than for CNF
or cutting planes, with the extra ingredient that the target formalism is a \emph{network} whose energy
we identify with the objective.

\medskip\noindent\textbf{Formalizing this problem, and problems like it.} Malekpour verified, in
PVS, a closed-form algorithm producing an $n$-queens solution for every $n > 3$~\citeref{malekpour2020}.
That is prior formal verification of $n$-queens, and the contrast is sharp: it verifies a
\emph{construction} for the variant that is always solvable, whereas we verify a \emph{reduction}
for the variant that is \textsf{NP}-complete, and it is the reduction that makes negative answers
meaningful. Gäher and Kunze mechanised the Cook--Levin theorem in Coq~\citeref{gaeher2021},
formalizing polynomial-time reductions in general rather than one concrete encoding for one solver.
Closest in shape to this paper, and at this venue, is Weber's SAT-based Sudoku
solver~\citeref{weber2006}: a puzzle encoded into a solver formalism and evaluated
empirically, without a proof that the encoding is correct.

\section{What we do not claim}\label{sec:notclaim}

\begin{itemize}\itemsep2pt
\item No claim of speed. See \secref{sec:meas}.
\item No new algorithm, and no improvement on the neural-network literature for this problem. The
search is~\citeref{liwang2022} unchanged. Ma\'{n}dziuk~\citeref{mandziuk1995} solved $n$-queens thirty
years ago with a binary Hopfield-type network over the same $\{0,1\}$ state space, using $N^2$
neurons and no auxiliary variables, and reports a convergence rate of $100\%$ (once $99\%$) on
boards of size $n = 4,8,16,32,48,64,80$, with the average number of outer iterations fitted by
$2.82\,n^{0.67}$. On the boards where a comparison is meaningful --- the empty ones --- our
population search is comparable to that up to $N = 8$, and it then falls away: we fail at $N = 10$
where he succeeds at $N = 80$.

We can point to three reasons, and none of them favours us. His update is a Hopfield update and
therefore descends the energy at every step, which is why it converges; the recurrences we inherit
from~\citeref{liwang2022} accumulate their input and have no descent property, as
\secref{sec:net} explains. His encoding is built for this problem and needs $N^2$ neurons, whereas a
generic reduction into $\lVert Ax-b\rVert^2$ costs us $N^2+4N-6$, the extra ones being slacks that
carry no information but must still be got right. And the constants of \secref{sec:solver} were
recovered on Sudoku, where the search is preceded by a propagation stage that has no analogue here.

What is new in this paper is therefore the proof and not the performance, and a reader wanting an
effective neurodynamic $n$-queens solver should use~\citeref{mandziuk1995} rather than this.
\item The search is a heuristic, and we prove nothing about whether it succeeds. We prove that its
positive answers are correct, and that its objective has a zero precisely when the board has a
completion.
\item For the backtracking baseline we prove soundness only, as explained in \secref{sec:meas}.
\item Nothing about cooling schedules. \secref{sec:net} concerns the equilibrium distribution at a
fixed temperature, in the limit as that temperature falls. It does not say how slowly a real run
must be cooled in order to stay near equilibrium. This distinction matters here precisely because
our single-machine measurements are poor: the theorem says the difficulty lies in the schedule and
not in the encoding, and it must not be read as a remedy.
\item Nothing about the accumulating recurrences of eqs.~(3) and (6), as explained in
\secref{sec:net}.
\end{itemize}

\section*{Acknowledgements}
% TODO: acknowledgements
\bibliographystyle{plain}
\begin{thebibliography}{99}

\bibitem{hnbm}
M.~Cipollina, M.~Karatarakis, and F.~Wiedijk.
\newblock Formalized Hopfield networks and Boltzmann machines.
\newblock In \emph{LPAR-26}, EPiC Series in Computing, 2026. To appear.

\bibitem{bell2009}
J.~Bell and B.~Stevens.
\newblock A survey of known results and research areas for $n$-queens.
\newblock \emph{Discrete Mathematics}, 309(1):1--31, 2009.

\bibitem{gent2017}
I.~P.~Gent, C.~Jefferson, and P.~Nightingale.
\newblock Complexity of $n$-Queens Completion.
\newblock \emph{Journal of Artificial Intelligence Research}, 59:815--848, 2017.

\bibitem{glock2022}
S.~Glock, D.~Munh\'{a} Correia, and B.~Sudakov.
\newblock The $n$-queens completion problem.
\newblock \emph{Research in the Mathematical Sciences}, 9:41, 2022. arXiv:2111.11402.

\bibitem{simkin2021}
M.~Simkin.
\newblock The number of $n$-queens configurations.
\newblock arXiv:2107.13460, 2021.

\bibitem{mandziuk1995}
J.~Ma\'{n}dziuk.
\newblock Solving the $N$-Queens problem with a binary Hopfield-type network.
\newblock \emph{Biological Cybernetics}, 72:439--445, 1995. DOI 10.1007/BF00201419.

\bibitem{merida2001}
E.~M\'{e}rida, J.~Mu\~{n}oz, and R.~Ben\'{i}tez.
\newblock A recurrent multivalued neural network for the $N$-Queens problem.
\newblock In \emph{Connectionist Models of Neurons, Learning Processes and Artificial
Intelligence}, volume 2084 of LNCS. Springer, 2001.

\bibitem{che2019}
H.~Che and J.~Wang.
\newblock A collaborative neurodynamic approach to global and combinatorial optimization.
\newblock \emph{Neural Networks}, 114:15--27, 2019. DOI 10.1016/j.neunet.2019.02.002.

\bibitem{liwang2022}
H.~Li and J.~Wang.
\newblock Collaborative neurodynamic algorithms for solving Sudoku puzzles.
\newblock In \emph{12th International Conference on Information Science and Technology (ICIST)},
pages 8--17. IEEE, 2022.

\bibitem{lucas2014ising}
A.~Lucas.
\newblock Ising formulations of many NP problems.
\newblock \emph{Frontiers in Physics}, 2:5, 2014. arXiv:1302.5843.

\bibitem{glover2019}
F.~Glover, G.~Kochenberger, and Y.~Du.
\newblock Quantum bridge analytics I: a tutorial on formulating and using QUBO models.
\newblock arXiv:1811.11538, 2019.

\bibitem{codognet2023}
P.~Codognet.
\newblock Encoding the at-most-one constraint for QUBO and quantum annealing: experiments with the
$N$-Queens problem.
\newblock In \emph{Companion Proceedings of the Genetic and Evolutionary Computation Conference
(GECCO '23)}. ACM, 2023. DOI 10.1145/3583133.3596394.

\bibitem{codel2023}
C.~R.~Codel, J.~Avigad, and M.~J.~H.~Heule.
\newblock Verified encodings for SAT solvers.
\newblock In \emph{Formal Methods in Computer-Aided Design (FMCAD 2023)}. TU Wien Academic Press,
2023.

\bibitem{szeider2026}
S.~Szeider.
\newblock PBLean: pseudo-Boolean proof certificates for Lean~4.
\newblock arXiv:2602.08692, 2026.

\bibitem{bogaerts2022}
B.~Bogaerts, S.~Gocht, C.~McCreesh, and J.~Nordstr\"{o}m.
\newblock Certified symmetry and dominance breaking for combinatorial optimisation.
\newblock arXiv:2203.12275, 2022.

\bibitem{subercaseaux2023}
B.~Subercaseaux and M.~J.~H.~Heule.
\newblock The packing chromatic number of the infinite square grid is 15.
\newblock In \emph{TACAS 2023}. Springer, 2023. arXiv:2301.09757.

\bibitem{malekpour2020}
M.~R.~Malekpour.
\newblock Formal verification of a solution to the $n$-Queens problem.
\newblock Technical Memorandum NASA/TM-2020, NASA Langley Research Center, April 2020.

\bibitem{gaeher2021}
L.~G\"{a}her and F.~Kunze.
\newblock Mechanising complexity theory: the Cook--Levin theorem in Coq.
\newblock In \emph{Interactive Theorem Proving (ITP 2021)}, volume 193 of LIPIcs, article 20, 2021.

\bibitem{weber2006}
T.~Weber.
\newblock A SAT-based Sudoku solver.
\newblock In \emph{LPAR-12 Short Paper Proceedings}, pages 11--15, 2006.

\end{thebibliography}

\end{document}
